\scanpage{147}
on voit que si $\lambda=(\lambda_\alpha)\in\unsure{O_k}(A)^I$, $(x_\alpha^*)\in\left(\prod_{\alpha\in I}V(N_\alpha)\right)(A)$
\[
u^*((\lambda_\alpha x_\alpha^*)_{\alpha\in I})=(u^\lambda)^*(x_\alpha^*)
\]
i.e.
\[
u^*\prod_\alpha\lambda_{\alpha\,V(N_\alpha)}=(u^\lambda)^*\ ,
\]
où \deleted{$u^\lambda$ est}
\[
u^\lambda=\sum_{(i_\alpha)}\left(\prod_\alpha\lambda_\alpha^{i_\alpha}\right)u_{(i_\alpha)}
\]
Donc, si on se donne $\underline n=(n_\alpha)_{\alpha\in I}\in\mathbb{N}^I$, $u^*$ est multihomogène de \inserted{multi}degré $\underline n$ ssi on a
\[
u:M\longrightarrow\bigotimes_\alpha\operatorname{Sym}^{n_\alpha}(N_\alpha)
\]
En particulier, si tous les $n_\alpha$ sont égaux à 1, cela montre que $u^*$ est hom. de degré 1 en tous les arguments ssi
\[
u:M\longrightarrow\bigotimes_\alpha N_\alpha
\]
(en identifiant $N_\alpha$ à $\operatorname{Sym}^1N_\alpha$). Mais alors on constate tout de suite que $u^*$ est \inserted{$=$} multilinéaire (\unsure{pouvant} aussi se déduire formellement du cas d'1 variable (``en fixant les autres'')). Donc les accouplements bil.
\[
\prod V(N_i)\longrightarrow\deleted{\unsure{V(N}}M
\]
correspondent aux hom
\[
M\longrightarrow\bigotimes_\alpha N_\alpha
\]
en particulier les accouplements bil.
\[
V(N)\times V(N')\mathbin{\text{---}}V(M)
\]
correspondent aux hom
\[
M\mathbin{\text{---}}N\otimes N'.
\]
\margin{146}
